\chapter{Functions as Values}\label{ch-function-values}

A function in Zeno is a value, like a number or a list. It can be
given a name, passed to another function, kept in a list or a dict,
and given back by a function. This chapter is about what that makes
possible.

\section{Passing a function}

A function's name, without brackets after it, is the function itself.
Passing it hands the function over, to be called later:

\begin{zeno}
function twice(f: function(int) -> int, x: int) -> int {
    return f(f(x))
}

function add_one(n: int) -> int {
    return n + 1
}

print(twice(add_one, 5))
print(add_one)
\end{zeno}

\verb|twice| calls whatever function it was given, twice. You have
already seen this with \verb|map|, \verb|filter| and \verb|reduce| in
\hyperref[ch-lists]{Lists}.

\section{The type of a function}

The type of a function is written as its heading with the names and
the body taken out:

\begin{center}
\begin{tabular}{ll}
\toprule
Function & Its type \\
\midrule
\verb|function add_one(n: int) -> int| & \verb|function(int) -> int| \\
\verb|function add(a: int, b: int) -> int| & \verb|function(int, int) -> int| \\
\verb|function report(text: str)| & \verb|function(str)| \\
\verb|function now() -> int| & \verb|function() -> int| \\
\bottomrule
\end{tabular}
\end{center}

A function fits a type when it can be called with those arguments and
gives back what the type says. One that does not fit is refused at the
call, with the reason:

\begin{zeno}[fails=yes]
function twice(f: function(int) -> int, x: int) -> int {
    return f(f(x))
}
function shout(text: str) -> str {
    return text.upper()
}
print(twice(shout, 1))
\end{zeno}

\begin{why}
There is nothing new to learn: the same keyword, the same brackets,
the same \verb|->| as a heading. A form such as \verb|(int) -> int|,
as in TypeScript and Swift, would be shorter, but it gives brackets a
second meaning in types. (Decision 1.11a.)
\end{why}

\section{A function with no name}

A function that is used only once can be written where it is used,
with the name left out:

\begin{zeno}
let double = function(n: int) -> int { return n * 2 }
print(double(21))

let readings = [15.0, 25.0, 30.0]
print(readings.map(function(c) { return c * 1.8 + 32.0 }))
print(readings.filter(function(c) { return c > 20.0 }))
\end{zeno}

When such a function is passed where a function type is expected, as
with \verb|map| here, its types can be left out: Zeno works them out
from the type it is given.

\begin{why}
It is exactly a function declaration with the name left out, so there
is only one way to write a function. Python's \verb|lambda n: n * 2|
is a second, smaller syntax that cannot hold statements, and
JavaScript's \verb|n => n * 2| would need the arrow that Zeno left out
of \verb|match|. (Decision 3.6a.)
\end{why}

For the commonest case, making a new list from an old one, a
comprehension is shorter still: \verb|[c * 1.8 + 32.0 for c in readings]|.

\section{Functions in lists and dicts}

Because a function is a value, it can be kept anywhere a value can. A
dict from names to functions is a simple way to run a command by name:

\begin{zeno}
function add(a: int, b: int) -> int {
    return a + b
}
function multiply(a: int, b: int) -> int {
    return a * b
}

let operations: dict[str, function(int, int) -> int] = [:]
operations["add"] = add
operations["times"] = multiply
operations["minus"] = function(a: int, b: int) -> int { return a - b }

for name in ["add", "times", "minus"] {
    print(name, operations[name](6, 3))
}
\end{zeno}

\section{Functions that make functions}

A function can give back a function. The function it gives back can
use the variables of the function that made it, even after that
function has returned:

\begin{zeno}
function scaler(factor: float) -> function(float) -> float {
    return function(x: float) -> float { return x * factor }
}

let to_millivolts = scaler(1000.0)
let halve = scaler(0.5)
print(to_millivolts(3.3), halve(10.0))
\end{zeno}

\verb|to_millivolts| and \verb|halve| are two functions made by one
\verb|scaler|, each with its own \verb|factor|. A function together
with the variables it keeps is called a \emph{closure}.

\subsection{It keeps the variable, not a copy}

A closure keeps the variable itself. If it changes the variable, the
change lasts until it is called again:

\begin{zeno}
function counter() -> function() -> int {
    var count = 0
    return function() -> int {
        count += 1
        return count
    }
}

let next_id = counter()
print(next_id(), next_id(), next_id())

let other = counter()
print(other())
\end{zeno}

Each call of \verb|counter| makes a new \verb|count|, so \verb|other|
starts again from 1.

\begin{why}
This is what the code says when it is read: there is one
\verb|count|. If the inner function had a copy, \verb|next_id()|
would give 1 every time, and nothing would say why. (Decision 3.6c.)
\end{why}

\subsection{Each turn of a loop has its own name}

The name in a \verb|for| loop is new on each turn, so a function made
inside a loop keeps the value of that turn:

\begin{zeno}
var handlers: list[function()] = []
for at in range(3) {
    handlers.append(function() { print("handler", at) })
}
for handle in handlers {
    handle()
}
\end{zeno}

In JavaScript with \verb|var|, and in Python, all three would print the
last value, 2: a famous puzzle that Zeno does not have.

\subsection{What it can see is fixed where it is written}

A function can see the variables declared above it, in the blocks
around it, and never one declared further down, even one with the
same name:

\begin{zeno}
let level = 1
function outer() -> int {
    let read = function() -> int { return level }
    let level = 2
    return read()
}
print(outer())
\end{zeno}

\verb|read| was written where the only \verb|level| in sight was the
one at the top, so that is the one it uses, and \verb|outer| gives 1.

\begin{why}
The meaning of a name should be settled by where it is written, not
by what happens to be declared later. This rule, called
\emph{lexical scope}, is the one nearly every modern language follows,
and it lets the compiler decide once which variables a function must
take with it. (Decisions 3.6c and 3.6d.)
\end{why}

\section{A function that raises}

A function that can raise says so in its heading, and its type says so
too, after what it gives back:

\begin{zeno}
function parse(text: str) -> int raises ValueError {
    if text == "x" {
        raise ValueError("not a number: x")
    }
    return int(text)
}

function total(texts: list[str], convert: function(str) -> int raises ValueError) -> int raises ValueError {
    var sum = 0
    for text in texts {
        sum += convert(text)
    }
    return sum
}

print(total(["1", "2", "3"], parse))
try {
    print(total(["4", "x"], parse))
} except ValueError as e {
    print("caught:", e)
}
\end{zeno}

A function fits \verb|raises ValueError| if \verb|except ValueError|
would catch everything it raises, so one that raises nothing fits as
well. Where the type says nothing, \verb|zeno run| and an ordinary
build let a raising function in, as they always did. A firmware build
holds it, because there the error travels as a value that each call
through the type has to test:

\begin{zeno}[file=apply.zeno, cmd=zeno build --firmware {file}, fails=yes]
function parse(text: str) -> int raises ValueError {
    return int(text)
}
function apply(f: function(str) -> int, text: str) -> int {
    return f(text)
}
print(apply(parse, "3"))
\end{zeno}

\begin{why}
A call through a function value is checked as a call of the function
is: \verb|total| must say that it raises \verb|ValueError| or catch it.
That works only if the type carries the kind of error, as Java's
\verb|throws| does on a method. It is also what a firmware build needs
to know to test every call through the value. (Decisions 56.1 to
56.5; new in 0.8.)
\end{why}

\section{A method as a value}

A method written without the brackets of a call is a function, with
its object already in it:

\begin{zeno}
class Counter {
    count: int
    function bump(self, by: int) -> int {
        self.count += by
        return self.count
    }
}

let c = Counter(0)
let step = c.bump
print(step(5), step(1), [1, 2].map(c.bump), c.count)
\end{zeno}

Its type is the method's, without \verb|self|: \verb|c.bump| is a
\verb|function(int) -> int|. It holds the object that \verb|c| held on
that line, so a \verb|var| given another object later does not change
it. (Decision 56.9; it builds since 0.8.)

\section{Calling something that is not a function}

Only a function can be called, and the checker knows which names hold
one:

\begin{zeno}[fails=yes]
let speed = 5
speed()
\end{zeno}

\begin{summary}
A function is a value: its name without brackets is the function. Its
type is written \verb|function(int, str) -> bool|. A function with no
name is written \verb|function(n) { return n * 2 }|. Functions can be
passed, stored in lists and dicts, and given back by other functions.
A type can say what its function raises, and a method without its
brackets is a function with its object in it.
A function given back keeps the variables around it, the variables
themselves and not copies, as they were in sight where it was written.
\end{summary}

\begin{tryit}
\begin{enumerate}
\item Write a function \verb|apply_all| that takes
  \verb|fs: list[function(int) -> int]| and \verb|x: int|, and gives
  each function's answer for \verb|x| in a \verb|list[int]|.
\item Write a function \verb|compose| that takes two functions of type
  \verb|function(int) -> int|, \verb|f| and \verb|g|, and gives a
  function of that type which does \verb|g| and then \verb|f|.
\item Make a counter that starts from a number and counts up by a step,
  both given to the function that makes it.
\item Build a small calculator: read lines such as \verb|add 2 3| with
  \verb|input()|, split them, and look the word up in a dict of
  functions. What should happen with a word that is not in the dict?
\end{enumerate}
\end{tryit}
